\documentclass[11pt,a4paper]{article}
\usepackage[T1]{fontenc}
\usepackage{lmodern,amsmath,amssymb,amsthm,mathtools,booktabs,microtype}
\usepackage[margin=27mm]{geometry}
\usepackage[colorlinks=true,linkcolor=blue,citecolor=blue,urlcolor=blue]{hyperref}
\setlength{\emergencystretch}{2em}
\newtheorem{theorem}{Theorem}[section]
\newtheorem{lemma}[theorem]{Lemma}
\newtheorem{proposition}[theorem]{Proposition}
\newtheorem{corollary}[theorem]{Corollary}
\newcommand{\Q}{\mathbb Q}\newcommand{\Z}{\mathbb Z}
\DeclareMathOperator{\Tr}{Tr}\DeclareMathOperator{\Gal}{Gal}
\title{An obstruction to finite M\"obius coupling of six elliptic contacts}
\author{}\date{17 November 2026}
\begin{document}\maketitle
\begin{abstract}
The six-contact polynomial of the conditional elliptic source has
exactly twelve rational M\"obius symmetries. Coupling its two arguments
by these symmetries fails the prescribed-index completeness assertion
for every fixed integer power at least four, and for every fixed finite
menu of such powers. At input one, ten branches have no rational slope
and the identity violates a contact guard. For the remaining involution,
an arithmetic progression of coefficient shifts produces a simple
contact at an inert prime, forcing a negative odd detector pole.
A binomial root-field lemma and Chebotarev's theorem supply the required
primes outside every fixed finite exceptional set. Distinct primes and
the Chinese remainder theorem give one progression obstructing all
powers in a finite menu simultaneously.
\end{abstract}

\section{The prescribed source}

Let
\[
 K=\Q(\sqrt2),\qquad D=(3-\sqrt2)^2,\qquad E:y^2=x^3-Dx,
\]
\[
 P=\left(\frac{57-28\sqrt2}{8},
                    \frac{-660+417\sqrt2}{32}\right).
\]
Conjugation in $K/\Q$ is denoted by a bar. For points $A,H,U$, use
$f=1/x$, $z=-x/y$ and the six tests
\begin{equation}\label{eq:tests}
\begin{gathered}
 r=\frac{z(U)}{z(H)},\qquad s=\frac{z(U-A)}{z(H-P)},\\
 C_1=f(H),\quad C_2=f(U),\quad
 C_3=\overline{f(H-P)},\quad C_4=\overline{f(U-A)},\\
 C_5=r-\bar s,\qquad C_6=r+\bar s-2X.
\end{gathered}
\end{equation}
All point, difference and quotient guards are retained. For an odd
integer $B\ge3$ and an integer $e\ge4$, define
\begin{equation}\label{eq:source}
 T=\Tr x(A),\qquad a_e=(X^2+B)^e+T^e,
 \qquad t_0=\lambda,\quad t_1=a_e\lambda,
\end{equation}
and
\begin{equation}\label{eq:F}
 F(t)=\frac{t(t+1)(2t+1)(3t+1)(4t+1)(5t+2)}{120}.
\end{equation}
For nonzero rational multipliers $\varepsilon_j$, the detectors are
\begin{equation}\label{eq:detector}
 d_j=\varepsilon_jF(t_j),\quad
 w_j=\frac{d_j}{4(t_j-1)^8},\quad
 W_j=\prod_{i=1}^6\left(1+\frac{C_i^2}{d_j}\right),\quad
 V_j=\Tr(W_j^{-1}),\quad b_j=\frac{4w_j}{V_j^2}.
\end{equation}
The literal source has $\varepsilon_0=\varepsilon_1=1$
\cite{Reweighted}. Its guards include
$d_j(t_j-1)V_j\ne0$ and invertibility of $W_j$.
The pole predicates considered here imply
\begin{equation}\label{eq:pole}
 v_p(b_j)\text{ is not a negative odd integer}\qquad(p\notin S),
\end{equation}
for a common fixed finite set $S$ of rational primes. An intrinsic
pole block for $7b_j$ has this property\cite{Intrinsic}.

Put
\[
 \mathcal R=\{0,-1,-1/2,-1/3,-1/4,-2/5\},\qquad
 \mathcal M=\{\phi\in\operatorname{PGL}_2(\Q):\phi(\mathcal R)=\mathcal R\}.
\]
The additional coupling is
\begin{equation}\label{eq:coupling}
                         a_e\lambda=\phi(\lambda),\qquad\phi\in\mathcal M,
\end{equation}
with a nonzero M\"obius denominator.

\begin{theorem}\label{thm:main}
Let $\mathcal E$ be any nonempty finite set of integers at least
four. Let $S$ be a finite set of primes containing $2,3,5,7$, and let
$\mathcal U\subset\Q^\times$ be finite and nonempty.
There are an odd integer $B_0$ and a positive even integer $M$ such
that, for every $B\ge3$ with $B\equiv B_0\pmod M$, the prescribed data
\[
                             X=1,\qquad A=P
\]
admit no guarded rational auxiliary tuple satisfying
\eqref{eq:source}--\eqref{eq:coupling} and at least one of the two
pole conditions~\eqref{eq:pole}, for any choices
\[
                  e\in\mathcal E,\qquad \phi\in\mathcal M,
                  \qquad \varepsilon_0,\varepsilon_1\in\mathcal U.
\]
The exponent, symmetry, multipliers and retained pole index may be
chosen separately for each tuple and shift. Pole predicates may vary
within any family whose members all imply~\eqref{eq:pole} outside
the same $S$.
\end{theorem}

Taking $\mathcal E=\{e\}$ gives the obstruction for each fixed integer
exponent $e\ge4$. The theorem excludes eventual prescribed-index
completeness for every fixed finite power menu. It concerns the
displayed scales and the prescribed point $A=P$: it does not assert
failure for every individual shift, for an infinite menu of exponents,
for an independently chosen root index, or for a different contact
system. If a restricted source retains the conclusion that the root
index equals $X$, the theorem also obstructs that source's eventual
completeness.

\section{The symmetry menu and uniform square gaps}

\begin{lemma}\label{lem:menu}
The group $\mathcal M$ consists of the following twelve maps
$\phi(t)=(\alpha t+\beta)/(\gamma t+\delta)$.
The last column gives the multiplier in
$F(\phi(t))(\gamma t+\delta)^6=\kappa F(t)$.
\begin{center}
\begin{tabular}{rrrrr}\toprule
$\alpha$&$\beta$&$\gamma$&$\delta$&$\kappa$\\\midrule
1&0&0&1&1\\
1&0&$-5$&$-1$&$-1$\\
1&1&2&$-1$&27\\
4&1&$-7$&$-1$&$-27$\\
2&1&$-3$&$-2$&$-1$\\
3&1&$-7$&$-2$&1\\
2&1&$-7$&$-3$&1\\
3&1&$-8$&$-3$&$-1$\\
1&1&$-7$&$-4$&$-27$\\
4&1&$-13$&$-4$&27\\
5&2&$-11$&$-5$&27\\
5&2&$-14$&$-5$&$-27$\\\bottomrule
\end{tabular}
\end{center}
\end{lemma}
\begin{proof}
A M\"obius map is determined by the images of $0,-1,-1/2$.
Enumerate their $6\cdot5\cdot4=120$ ordered triples of distinct
images in $\mathcal R$. For an image triple $(u,v,w)$, a representing
matrix is
\[
 \delta=v-w,\qquad \gamma=u+v-2w,\qquad
 \alpha=v\gamma+(u-v)\delta,\qquad \beta=u\delta.
\]
Retain the nonsingular maps having no pole in $\mathcal R$ and
mapping all six roots onto $\mathcal R$. Normalization gives exactly
the displayed matrices. Multiplying the six linear factors in
\eqref{eq:F} verifies the multipliers. These finite rational
calculations are reproduced by the standard-library checker.
\end{proof}

At the prescribed point, $T=57/4$. Write $a=a_e=p_B/q$, where
\begin{equation}\label{eq:scale}
              q=4^e,\qquad p_B=q(B+1)^e+57^e.
\end{equation}
In particular $a>4$. The identity symmetry forces $\lambda=0$,
contrary to the contact guard. In each of rows $3$--$12$, a rational
solution of the coupling would make
\[
                  G=(\delta p_B-\alpha q)^2+4\gamma\beta p_Bq
\]
an integer square. Set
\begin{equation}\label{eq:gapdefinitions}
 c=-\delta\alpha+2\gamma\beta,\qquad
 N=\delta^2p_B+cq,\qquad
 K_0=4\beta\gamma(\alpha\delta-\beta\gamma).
\end{equation}
Expansion gives $\delta^2G=N^2+K_0q^2$.

\begin{lemma}\label{lem:gap}
For every integer $e\ge4$, every integer $B\ge3$, and all rows
$3$--$12$, one has $N>0$ and
\[
                         0<|K_0|q^2<2N-1.
\]
Consequently none of these ten branches has a rational slope.
\end{lemma}
\begin{proof}
The integer $N$ increases with $B$, so take $B=3$. The gap is
\[
 g_e=2\delta^2 57^e+(2\delta^2-|K_0|)16^e+2c4^e-1.
\]
The initial values are positive:
\begin{center}
\begin{tabular}{rrrr}\toprule
Row&$c$&$K_0$&$g_4$\\\midrule
3&5&$-24$&19672769\\
4&$-10$&$-84$&15732929\\
5&$-2$&12&84184839\\
6&$-8$&$-28$&83133191\\
7&$-8$&$-28$&189348561\\
8&$-7$&32&189086929\\
9&$-10$&$-84$&334379039\\
10&$-10$&156&329660447\\
11&$-19$&264&513765617\\
12&$-31$&$-336$&509040881\\\bottomrule
\end{tabular}
\end{center}
Moreover
\[
        g_{e+1}-16g_e=82\delta^2 57^e-24c4^e+15>0,
\]
because $\delta^2\ge1$, $c\le5$, and $57^e>2\cdot4^e$.
Thus $g_e>0$ for all integers $e\ge4$. This also gives $N>0$.
If $K_0<0$, the integer $N^2+K_0q^2$ lies strictly between
$(N-1)^2$ and $N^2$; if $K_0>0$, it lies strictly between
$N^2$ and $(N+1)^2$. It cannot be a square.
\end{proof}

\section{Simple contacts at inert primes}

The remaining symmetry is $\phi(t)=-t/(5t+1)$. Since $\lambda=0$
is forbidden, its coupling has the unique candidate
\begin{equation}\label{eq:involution}
 \lambda=-\frac{a+1}{5a},\qquad
 H(a)=(a^2-1)(4a-1)(3a-2)(2a-3)(a-4).
\end{equation}
Direct multiplication gives
\begin{equation}\label{eq:H}
             -F(\lambda)=\frac{H(a)}{375000a^6},\qquad
              F(a\lambda)=\frac{H(a)}{375000},\qquad H'(4)=11250.
\end{equation}

\begin{lemma}\label{lem:inert}
Let $p\ne2,7$ be inert in $K$. For a defined detector
\[
 W=\prod_i(1+C_i^2/d),\qquad V=\Tr(W^{-1}),\qquad b=4w/V^2,
\]
suppose $v_p(d)=v_p(w)=1$ and the tests include
$f(H)$ and $\overline{f(H-P)}$. Then either $V=0$ or $v_p(b)$
is a negative odd integer.
\end{lemma}
\begin{proof}
The curve has good reduction at $p$, since its discriminant is
$64D^3$ and $\operatorname{N}_{K/\Q}(D)=49$. The coordinates of $P$
are integral at $p$, so its reduction is affine and not the identity.
Conjugation fixes the unique place above $p$. If both displayed tests
had positive valuation, both $H$ and $H-P$ would reduce to the identity,
which would force $P$ to do so. At least one test therefore has
nonpositive valuation.

In the unramified quadratic completion, valuations are integral and
$v(p)=1$. For a nonzero test, the two summands of $1+C_i^2/d$ have
different valuations, so its valuation is $\min(0,2v(C_i)-1)\le0$.
A zero test gives the factor one. At least one factor has valuation
at most $-1$. Hence $W^{-1}\in p\mathcal O_{K_p}$ and
$V\in p\Z_p$. If $V\ne0$, then $v_p(b)=1-2v_p(V)$ is negative
and odd.
\end{proof}

\section{A binomial root-field lemma}

\begin{lemma}\label{lem:rootfield}
Let $m$ be a positive odd integer, and let $e$ be a positive integer.
If $e$ is odd or $m$ is nonsquare, then $Y^e+m$ has a root $\alpha$
such that $\sqrt2\notin\Q(\alpha)$.
\end{lemma}
\begin{proof}
If $e$ is odd, the polynomial has an irreducible factor of odd degree
over $\Q$. The field generated by a root of that factor has odd degree,
so cannot contain a quadratic subfield.

Now suppose that $e$ is even and $m$ is nonsquare.
Write $e=2^r s$ with $r\ge1$ and $s$ odd. Choose an odd prime $\ell$
for which $v_\ell(m)$ is odd. Let $\beta$ satisfy
$\beta^{2^r}=-m$ and put $L=\Q(\beta)$. At any prime of $L$
above $\ell$, normalize the valuation by $v(\ell)=1$. Then
\[
                       v(\beta)=\frac{v_\ell(m)}{2^r}.
\]
The denominator in lowest terms is $2^r$, so the ramification index
is divisible by $2^r$. Since $[L:\Q]\le2^r$, the degree equals
$2^r$, and the degree formula shows that $\ell$ has exactly one
prime in $L$, with ramification index $2^r$ and residue degree one.

The quadratic field $\Q(\sqrt2)$ is unramified at the odd prime
$\ell$. If $\ell$ splits there, an inclusion $\Q(\sqrt2)\subset L$
would give at least two primes of $L$ above $\ell$. If $\ell$ is
inert there, the residue degree in $L$ would be divisible by two.
Both conclusions contradict the preceding paragraph, so
$\sqrt2\notin L$.

Factor $Y^s-\beta$ over $L$. Its odd degree implies that at least
one irreducible factor has odd degree. Choose a root $\alpha$ of
such a factor. Then $\beta=\alpha^s$, so $L\subset\Q(\alpha)$
and $[\Q(\alpha):L]$ is odd. Consequently $\sqrt2$ cannot belong
to $\Q(\alpha)$, since otherwise the quadratic extension
$L(\sqrt2)/L$ would be a subextension. Finally
$\alpha^e=\beta^{2^r}=-m$.
\end{proof}

\begin{corollary}\label{cor:primes}
For every integer $e\ge4$, infinitely many rational primes inert in
$\Q(\sqrt2)$ admit a simple residue root of
\[
                         Y^e+57^e-4^{e+1}.
\]
These primes may be chosen outside any fixed finite set.
\end{corollary}
\begin{proof}
Put $m=57^e-4^{e+1}$, which is positive and odd. If $e$ is even,
put $k=e/2\ge2$ and $n=57^k$. Since
$4\cdot16^k<57^k<2\cdot57^k-1$, one has
\[
                              (n-1)^2<m<n^2.
\]
Thus $m$ is nonsquare when $e$ is even. For either parity,
Lemma~\ref{lem:rootfield} gives a root $\alpha$ of $Y^e+m$ with
$\sqrt2\notin L=\Q(\alpha)$.

Let $M$ be the splitting field of $Y^e+m$ with $\sqrt2$ adjoined.
The nontrivial automorphism of $L(\sqrt2)/L$ extends to an element
$\sigma\in\Gal(M/L)$. It fixes $\alpha$ and sends $\sqrt2$ to
$-\sqrt2$. Chebotarev's theorem\cite[Theorem 8.31]{Milne} gives
infinitely many unramified rational primes whose Frobenius conjugacy
class is that of $\sigma$. Their Frobenius is nontrivial on
$\Q(\sqrt2)$, so these primes are inert there. Its action on the
roots has a fixed point, so, after excluding primes dividing the
polynomial discriminant, $Y^e+m$ has a simple root modulo each
such prime. Removing any additional finite set preserves infinitude.
\end{proof}

\section{Simultaneous obstruction for a finite menu}

\begin{proof}[Proof of Theorem~\ref{thm:main}]
Enlarge $S$ by all numerator and denominator primes of elements of
$\mathcal U$. For each $e\in\mathcal E$, choose a distinct prime
$p_e$ supplied by Corollary~\ref{cor:primes}, outside $S$ and the
primes dividing $e(57^e-4^{e+1})$. Choose a simple root $y_e$ modulo
$p_e$. The residue
\[
                        B+1\equiv y_e/4\pmod{p_e}
\]
then gives $a_e(B)\equiv4\pmod{p_e}$. Its derivative
$e(B+1)^{e-1}$ is nonzero modulo $p_e$. Of the $p_e$ lifts of
this residue class modulo $p_e^2$, exactly one satisfies
$a_e(B)\equiv4\pmod{p_e^2}$. Choose any other lift $b_e$.
Since $H'(4)=11250$ is a unit at $p_e$, this gives
\begin{equation}\label{eq:exactcontact}
               v_{p_e}(H(a_e(B)))=1
               \qquad(B\equiv b_e\pmod{p_e^2}).
\end{equation}

By the Chinese remainder theorem, the simultaneous conditions
\[
             B\equiv1\pmod2,\qquad
             B\equiv b_e\pmod{p_e^2}\quad(e\in\mathcal E)
\]
define one arithmetic progression of odd integers, of modulus
$M=2\prod_{e\in\mathcal E}p_e^2$.
Take any member $B\ge3$ and any choice $e\in\mathcal E$.
Lemma~\ref{lem:gap} excludes ten symmetries, and the identity
violates the contact guard. The remaining involution has
$\lambda=-(a_e+1)/(5a_e)$. Modulo $p_e$, its values are
\[
                     a_e=4,\qquad \lambda=-1/4,
                     \qquad a_e\lambda=-1.
\]
The denominators in~\eqref{eq:H}, both $t_j-1$, and the M\"obius
denominator $5\lambda+1$ are units. Equations~\eqref{eq:H}
and~\eqref{eq:exactcontact} therefore give valuation one for both
$F(t_j)$ and both weights $w_j$. Each allowed multiplier is a
$p_e$-adic unit, so both $d_j$ also have valuation one.
Lemma~\ref{lem:inert} contradicts either retained pole condition.
This argument applies to every exponent in the menu at the same
shift $B$, completing the proof.
\end{proof}

\section{Exact certificates}

The standard-library checker \texttt{check.py} verifies all $120$
ordered image triples, the twelve multiplier identities, the ten
discriminant identities, the positive initial square gaps, the
coefficient inequalities in their recurrence, and the involution
identities. It also verifies a simultaneous local certificate for
the following finite menu. The residue $b_e$ is taken modulo $p_e^2$.
\begin{center}
\begin{tabular}{rrrr}\toprule
$e$&$p_e$&$y_e$&$b_e$\\\midrule
4&76213&11332&2832\\
5&76333&44199&30132\\
6&76259&34651&65856\\
7&76261&16102&42155\\
8&76243&36746&47307\\
10&76253&28091&26085\\
12&76379&10315&59862\\
16&76283&9710&40568\\
32&76421&24909&63542\\\bottomrule
\end{tabular}
\end{center}
Each prime exceeds $76176$ and is verified prime by trial division.
The checker verifies inertness, the simple binomial root, the exact
contact valuation modulo $p_e^2$, the relevant denominator units,
and the Chinese remainder progression. The infinite assertions for
arbitrary integer powers at least four and arbitrary finite menus follow from the
proofs above, independently of the numerical examples.

\begin{thebibliography}{9}
\bibitem{Reweighted}
\emph{Quadratic reweightings of native elliptic detectors}, repository source,
\nolinkurl{papers/pole_parameters/reweighted_elliptic_detectors_20261116_v53.tex}.
\bibitem{Intrinsic}
\emph{A two-coordinate pole test of coefficient height four}, repository source,
\nolinkurl{papers/pole_parameters/intrinsic_pole_20261117_v56.tex}.
\bibitem{Milne}
J.~S.~Milne, \emph{Algebraic Number Theory}, version 3.08, 2020,
Theorem~8.31, \url{https://www.jmilne.org/math/CourseNotes/ANT.pdf},
accessed 17 November 2026.
\end{thebibliography}
\end{document}
